% Part: proof-theory
% Chapter: cut-elimination
% Section: introduction

\documentclass[../../../include/open-logic-section]{subfiles}

\begin{document}

\olfileid{pt}{cut}{int}

\olsection{పరిచయం}

\Cut{} నియమం ఇది:
\[
\Axiom$\Gamma \fCenter \Delta, !A$
\Axiom$!A, \Pi \fCenter \Lambda$
\RightLabel{\Cut}
\BinaryInf$\Gamma, \Pi \fCenter \Delta, \Lambda$
\DisplayProof
\]
(ఇక్కడ \tetoken{సూత్రం}{!!{formula}}~$!A$ను \emph{కట్ సూత్రం} అంటారు.)

జెంట్సెన్ అసలు సీక్వెంట్ కలనం~\Log{LK}లో \Cut{} నియమం ఉంది. జెంట్సెన్ \emph{Hauptsatz} (“ప్రధాన ప్రమేయం”) ప్రకారం~\Log{LK}లో \tetoken{నిరూపణీయమైన}{!!{provable}} ప్రతి సీక్వెంట్ \Cut{} నియమం లేకుండానూ \tetoken{నిరూపణీయం}{!!{provable}}; అంటే~\Log{LK}-నిరూపణల నుంచి \Cut{}ను “తొలగించవచ్చు.” మన~\Log{G1c}, \Log{G3c} వ్యవస్థల్లో \Cut{} లేదు, అయినా వాటికీ ఇదే ప్రశ్న అడగవచ్చు: \Cut{}తో పాటు~\Log{G1c} (లేదా~\Log{G3c})లోని ఇతర నియమాలను ఉపయోగించే \tetoken{నిరూపణల}{!!{proof}s} నుంచి \Cut{}ను తొలగించవచ్చా? మన స్థిర పదజాలంలో సమాన ప్రశ్న: \Cut{} అనేది~\Log{G1c} (లేదా~\Log{G3c})లో అనుమతిత నియమమా?

కట్ తొలగింపు ప్రమేయం నిరూపణ సిద్ధాంతంలోని అత్యంత ముఖ్యమైన, ప్రసిద్ధమైన ఫలితాల్లో ఒకటి. మొదట, అవసరమని తొలిచూపుకు అనిపించే నియమం~\Log{G1c}, \Log{G3c}లకు అవసరం లేదని ఇది చూపుతుంది. నిజమైన నిరూపణలను రూపబద్ధం చేయడానికి ప్రయత్నిస్తే లెమ్మాల వాడకాన్ని నిర్వహించే మార్గం కావాలని తెలుస్తుంది. గణితజ్ఞులు ఒక ఫలితాన్ని నిరూపించి, తరువాత మరొక ఫలితం నిరూపణలో దాన్ని ఉపయోగిస్తారు. కాబట్టి ముందుగా ఫలితం~$L$ (లెమ్మా) నిరూపణను $\Sequent L$ సీక్వెంట్ యొక్క \tetoken{నిరూపణ}{!!{proof}}గా రూపబద్ధం చేయవచ్చు. తరువాత~$L$ను ఉపయోగించే రెండో ఫలితం~$R$ నిరూపణను $L \Sequent
R$ సీక్వెంట్ యొక్క \tetoken{నిరూపణ}{!!{proof}}గా రూపబద్ధం చేయవచ్చు. అప్పుడు~$R$కు నిరూపణ, అంటే కేవలం $\Sequent R$కు \tetoken{ఒక నిరూపణ}{!!a{proof}}, ఉందని ఎలా తెలుస్తుంది? ఆ రెండు \tetoken{నిరూపణలను}{!!{proof}s} కలిపే మార్గం కావాలి; \Cut{} నియమం దాన్ని ఇస్తుంది.

\Log{G1c}, \Log{G3c}లు సంపూర్ణమైనవని, అందువల్ల నిరూపించగలమని ఆశించే అన్ని సర్వసత్య సీక్వెంట్‌లనూ అవి నిరూపిస్తాయని చెప్పాం. దీన్ని నేరుగా నిరూపించవచ్చు. కానీ జెంట్సెన్ రాసిన కాలంలో స్వీకృత-ఆధారిత నిగమన వ్యవస్థల సంపూర్ణత మాత్రమే తెలిసింది. \Log{LK} సంపూర్ణత చూపడానికి జెంట్సెన్ స్వీకృత వ్యవస్థలోని \tetoken{నిగమనాలను}{!!{derivation}s} \Log{LK}-\tetoken{నిరూపణలుగా}{!!{proof}s} మార్చే అనువాదాన్ని ఇచ్చాడు. ఆ అనువాదంలో \Cut{} కీలకం. కట్ తొలగింపు సంప్రదాయ తర్కానికే ముఖ్యమైనది కాదు: అనేక ఇతర తర్కాలకు సీక్వెంట్ కలనాలు ఉన్నాయి. కొన్ని తర్కాల్లో ఆ కలనాల సంపూర్ణతను నేరుగా నిరూపించడం కష్టం లేదా అసాధ్యం; అప్పుడు ఇతర వ్యవస్థల నుంచి (\Cut{}ను వాడే) అనువాదాలే సంపూర్ణత చూపే మార్గం. తరువాత \Cut{} అనవసరమని, అది లేని వ్యవస్థ సంపూర్ణమని స్థాపించడానికి కట్ తొలగింపును నిరూపణ-సిద్ధాంత పద్ధతిలో నిరూపించాలి.

ఇతర స్వతంత్రంగా ఆసక్తికరమైన ఫలితాలను పొందడానికీ కట్ తొలగింపు ఉపయోగపడుతుంది. మనం పరిశీలించే రెండు ఉదాహరణలు క్రెయిగ్ అంతర్వర్తన లెమ్మా, హెర్బ్రాండ్ ప్రమేయం.

కట్ తొలగింపు నిరూపణలో సాధారణంగా \Cut{}ను ఉపయోగించే \tetoken{ఒక నిరూపణ}{!!a{proof}}ను, దాన్ని ఉపయోగించని నిరూపణగా ఎలా మార్చాలో చూపుతారు. ఈ మార్పులు సాధారణంగా “స్థానికం”: నిరూపణలో ఒక భాగాన్ని (ఎక్కువగా \Cut{} నిగమనంతో ముగిసే ఉప-\tetoken{నిరూపణ}{!!{proof}}ను) తీసుకుని, అదే సీక్వెంట్‌కు చెందిన మరొక ఉప-\tetoken{నిరూపణ}{!!{proof}}తో భర్తీ చేస్తాం; అందులో ఆ \Cut{} నిగమనాన్ని తొలగిస్తాం లేదా ఇతర నిగమనాలతో భర్తీ చేస్తాం. ఇటువంటి వాదనలు \Cut{} గల \tetoken{నిరూపణలను}{!!{proof}s} అది లేని నిరూపణలుగా దశలవారీగా మార్చే అల్గోరిథాలను ఇస్తాయి. అన్వయాల్లో కట్ తొలగింపు పద్ధతులు వాడినప్పుడు ఇది మరింత సమాచారం ఇస్తుంది. ఉదాహరణకు, మధ్యవర్తన లెమ్మాకు నమూనా-సిద్ధాంత నిరూపణలు “$!A \lif !B$ సర్వసత్యమైతే మధ్యవర్తి ఉంది” అనే రూపంలో ఉంటాయి (ఇప్పటికి మధ్యవర్తి అంటే ఏమిటనేది పక్కన పెడదాం). కట్ తొలగింపును ఉపయోగించే నిరూపణ-సిద్ధాంత వాదన $!A
\lif !B$ యొక్క \tetoken{ఒక నిరూపణ}{!!a{proof}} నుంచి మధ్యవర్తిని తిరిగి ఇచ్చే అల్గోరిథం ఇస్తుంది. నమూనా-సిద్ధాంత వాదనతో పోలిస్తే ఇది \emph{నిర్మాణాత్మకం}.

కట్ తొలగింపు ప్రమేయాన్ని నిరూపించే ప్రసిద్ధ మార్గాలు రెండు. మొదటిది (జెంట్సెన్‌ది) నిరూపణలో అత్యున్నత కట్‌లను తొలగించవచ్చని చూపి, కట్‌లేవీ మిగలకుండా అత్యున్నత కట్‌లను వరుసగా తొలగిస్తుంది. రెండోది (మొదట షుట్టె, టైట్ ప్రతిపాదించింది) \tetoken{ఒక నిరూపణ}{!!a{proof}}లో అత్యంత సంక్లిష్ట కట్‌లపై దృష్టి పెట్టి, మరింత లోతైన కట్ \tetoken{సూత్రం}{!!{formula}} గల మరొక కట్ లేని అర్థంలో గరిష్ఠమైన కట్‌లను వరుసగా తొలగిస్తుంది. రెండింటికీ ప్రయోజనాలు, పరిమితులు ఉన్నాయి. కొన్ని వ్యవస్థల్లో ఒక పద్ధతి పనిచేసినా మరొకటి కష్టమో అసాధ్యమో కావచ్చు. అందుకే రెండు పద్ధతులనూ వివరిస్తాం. ~\Log{G3c} కోసం కట్ తొలగింపును నిరూపిస్తాం. నిజానికి \Cut{}ను కాదు, \emph{సందర్భాన్ని పంచుకునే కట్}~\CutCS{}ను తొలగించవచ్చని చూపుతాం:
\[
\Axiom$\Gamma \fCenter \Delta, !A$
\Axiom$!A, \Gamma \fCenter \Delta$
\RightLabel{\CutCS}
\BinaryInf$\Gamma \fCenter \Delta$
\DisplayProof
\]
\sourcecorrection{OLTEPTCUTINT-001}{ఈ చిత్రం సందర్భాన్ని పంచుకునే \CutCS{} నియమాన్ని చూపుతుంది; మూలంలో లేబుల్‌ను పొరపాటున \Cut{}గా ఇచ్చారు.}
ఒక సీక్వెంట్ కలనంలో బలహీనీకరణ, సంకోచనం (వాస్తవ నియమాలుగా~\Log{G1c}లో ఉన్నా, అనుమతిత నియమాలుగా~\Log{G3c}లో ఉన్నా) సాధ్యమైతే, \Cut{} \CutCS{}ను అనుకరించగలదు; విపర్యయమూ సాధ్యం. మొదటి కట్ తొలగింపు పద్ధతిని~\CutCS{}తో వివరించడం సులభం, రెండో పద్ధతి~\CutCS{}తోనే పనిచేస్తుంది కనుక \Cut{}కు బదులు \CutCS{}ను ఎంచుకున్నాం.

\end{document}

